\section{\texorpdfstring{Пример с нетривиальной целевой функцией, в котором симплекс\=/метод возвращает точку, не являющуюся вершиной исходного квадрата}{Пример с нетривиальной целевой функцией, в котором симплекс-метод возвращает точку, не являющуюся вершиной исходного квадрата}}

Рассмотрим задачу линейного программирования
\begin{equation*}
  y \to \max
\end{equation*}
при ограничениях
\begin{equation*}
  -1 \leqslant x \leqslant 1,
\end{equation*}
\begin{equation*}
  0 \leqslant y \leqslant 1,
\end{equation*}
где переменная~$x$ считается свободной, а переменная~$y$ неотрицательна.

Допустимое множество в исходном пространстве переменных~$(x, y)$ представляет собой квадрат
\begin{equation*}
  Q = \big\{ (x, y) \in \mathbb{R}^2 \bigm| -1 \leqslant x \leqslant 1,\ 0 \leqslant y \leqslant 1 \big\}.
\end{equation*}
Его вершинами являются точки
\begin{equation*}
  (-1, 0), \qquad (1, 0), \qquad (-1, 1), \qquad (1, 1).
\end{equation*}

Целевая функция нетривиальна: её коэффициентный вектор в исходных переменных равен
\begin{equation*}
  c = (0, 1).
\end{equation*}
Максимум функции~$y$ на квадрате~$Q$ достигается на всём верхнем ребре
\begin{equation*}
  \big\{ (x, 1) \bigm| -1 \leqslant x \leqslant 1 \big\}.
\end{equation*}
Следовательно, точка
\begin{equation*}
  (0, 1)
\end{equation*}
является оптимальной, но не является вершиной квадрата~$Q$.

Покажем теперь, что после перехода к стандартной форме симплекс\=/метод может вернуть именно эту точку.

Поскольку переменная~$x$ свободна, представим её в виде разности двух неотрицательных переменных:
\begin{equation*}
  x = x^+ - x^-,
  \qquad
  x^+ \geqslant 0,
  \qquad
  x^- \geqslant 0.
\end{equation*}

Тогда задача приводится к стандартной форме
\begin{align}
  x^+ - x^- + s_1 &= 1,\nonumber\\
  -x^+ + x^- + s_2 &= 1,\nonumber\\
  y + s_3 &= 1,\nonumber\\
  x^+, x^-, y, s_1, s_2, s_3 &\geqslant 0,\label{temp:eq:standard-form}
\end{align}
где~$s_1, s_2, s_3 \geqslant 0$ --- дополнительные переменные, а целевая функция принимает вид
\begin{equation*}
  y \to \max.
\end{equation*}

В качестве начального базиса возьмём переменные~$s_1$, $s_2$, $s_3$.
Тогда небазисными будут переменные~$x^+$, $x^-$, $y$, и начальное базисное допустимое решение определяется условиями
\begin{equation*}
  x^+ = 0,
  \qquad
  x^- = 0,
  \qquad
  y = 0.
\end{equation*}
Подставляя эти значения в систему~\eqref{temp:eq:standard-form}, получаем
\begin{equation*}
  s_1 = 1,
  \qquad
  s_2 = 1,
  \qquad
  s_3 = 1.
\end{equation*}

Начальная симплекс\=/таблица имеет вид
\begin{equation*}
  \begin{array}{c|rrrrrr|r}
    \text{Базис} & x^+ & x^- & y & s_1 & s_2 & s_3 & b \\
    \hline
    s_1 & 1 & -1 & 0 & 1 & 0 & 0 & 1 \\
    s_2 & -1 & 1 & 0 & 0 & 1 & 0 & 1 \\
    s_3 & 0 & 0 & 1 & 0 & 0 & 1 & 1 \\
    \hline
    z & 0 & 0 & -1 & 0 & 0 & 0 & 0
  \end{array}
\end{equation*}
Последняя строка здесь соответствует канонической записи целевой функции $z - y = 0$.
Поскольку коэффициенты при базисных переменных в целевой функции равны нулю, строка оценок $C_j - Z_j$ равна
\begin{equation}
  \begin{array}{c|rrrrrr}
     & x^+ & x^- & y & s_1 & s_2 & s_3 \\
    \hline
    C_j - Z_j & 0 & 0 & 1 & 0 & 0 & 0
  \end{array}
  \label{temp:eq:reduced-costs-initial}
\end{equation}
Поэтому в базис следует вводить переменную~$y$, так как только она имеет положительную оценку в~\eqref{temp:eq:reduced-costs-initial}.

Рассмотрим столбец переменной~$y$.
Он имеет вид
\begin{equation*}
  \begin{pmatrix}
    0 \\
    0 \\
    1
  \end{pmatrix}.
\end{equation*}
Отношения $b_i / a_{ij}$ вычисляются только для положительных элементов этого столбца.
В первой и второй строках коэффициенты равны нулю, поэтому они не участвуют в выборе разрешающей строки.
В третьей строке имеем
\begin{equation*}
  \frac{1}{1} = 1.
\end{equation*}
Следовательно, из базиса выходит переменная~$s_3$, а разрешающим элементом является единица в третьей строке.

После pivot\=/перехода базисными становятся переменные~$s_1$, $s_2$, $y$, а новая таблица имеет вид
\begin{equation*}
  \begin{array}{c|rrrrrr|r}
    \text{Базис} & x^+ & x^- & y & s_1 & s_2 & s_3 & b \\
    \hline
    s_1 & 1 & -1 & 0 & 1 & 0 & 0 & 1 \\
    s_2 & -1 & 1 & 0 & 0 & 1 & 0 & 1 \\
    y & 0 & 0 & 1 & 0 & 0 & 1 & 1 \\
    \hline
    z & 0 & 0 & 0 & 0 & 0 & 1 & 1
  \end{array}
\end{equation*}
Здесь последняя строка соответствует канонической записи $z + s_3 = 1$.
В данном примере разрешающий элемент уже равен единице, а все остальные элементы в столбце~$y$ равны нулю, поэтому строки ограничений не изменяются численно; меняется только базис, строка цели и строка оценок.

Для нового базиса коэффициенты целевой функции равны $(0, 0, 1)$, поэтому строка оценок после pivot\=/перехода принимает вид
\begin{equation}
  \begin{array}{c|rrrrrr}
     & x^+ & x^- & y & s_1 & s_2 & s_3 \\
    \hline
    C_j - Z_j & 0 & 0 & 0 & 0 & 0 & -1
  \end{array}
  \label{temp:eq:reduced-costs-final}
\end{equation}
Поскольку в строке~\eqref{temp:eq:reduced-costs-final} все оценки неположительны, текущее базисное допустимое решение является оптимальным.

Теперь небазисными являются переменные~$x^+$, $x^-$, $s_3$.
Следовательно,
\begin{equation*}
  x^+ = 0,
  \qquad
  x^- = 0,
  \qquad
  s_3 = 0.
\end{equation*}
Из новой таблицы получаем
\begin{equation*}
  s_1 = 1,
  \qquad
  s_2 = 1,
  \qquad
  y = 1.
\end{equation*}

Итак, в расширенном пространстве найдено базисное допустимое решение
\begin{equation*}
  x^+ = 0,
  \qquad
  x^- = 0,
  \qquad
  y = 1,
  \qquad
  s_1 = 1,
  \qquad
  s_2 = 1,
  \qquad
  s_3 = 0.
\end{equation*}
Возвращаясь к исходной переменной, получаем
\begin{equation}
  x = x^+ - x^- = 0.
  \label{temp:eq:x-recovery}
\end{equation}
Следовательно, симплекс\=/метод возвращает в исходном пространстве точку
\begin{equation}
  (x, y) = (0, 1).
  \label{temp:eq:projected-solution}
\end{equation}
Значение целевой функции в точке~\eqref{temp:eq:projected-solution} равно
\begin{equation*}
  z = 1.
\end{equation*}

Остаётся заметить, что точка~$(0, 1)$ не является вершиной квадрата~$Q$.
Действительно,
\begin{equation*}
  (0, 1) = \frac{1}{2}(-1, 1) + \frac{1}{2}(1, 1),
\end{equation*}
причём точки~$(-1, 1)$ и~$(1, 1)$ различны и принадлежат множеству~$Q$.
Следовательно, точка~$(0, 1)$ лежит на верхнем ребре квадрата и не является его вершиной.

Тем самым получен двумерный пример с нетривиальной целевой функцией, в котором симплекс\=/метод находит базисное допустимое решение расширенной задачи, а его проекция~\eqref{temp:eq:projected-solution} на исходное пространство переменных не является вершиной исходного многоугольника.

Причина этого эффекта состоит в том, что переход
\begin{equation*}
  x = x^+ - x^-
\end{equation*}
не является взаимно однозначным.
Симплекс\=/метод действительно работает с вершинами многогранника в пространстве переменных
\begin{equation*}
  (x^+, x^-, y, s_1, s_2, s_3),
\end{equation*}
однако после перехода к исходным переменным вершина расширенного многогранника может перейти в точку, не являющуюся вершиной исходного квадрата.
